\documentclass[10pt,a4paper]{article}

\usepackage[utf8]{inputenc}
\usepackage[cyr]{aeguill}
\usepackage[francais]{babel}

% Les packages
%\usepackage[french]{babel}
\usepackage{fancybox}
\usepackage{graphics}
\usepackage{color}
\usepackage{latexsym}
\usepackage{amsmath,amsthm,amssymb}
\usepackage[plain]{fullpage} 
\usepackage{verbatim}
\usepackage{times,epsfig}
\usepackage{listings}
\usepackage{multirow}
\usepackage{amsmath,amsthm,amssymb,amscd,txfonts,bbm} % RAJOUT ICI !!!
\usepackage{graphics,epsfig} % RAJOUT ICI !!!
%\usepackage{algorithmic}
\usepackage[final]{pdfpages} 
\usepackage{fancyvrb}
\usepackage{enumerate}
\usepackage{listingsutf8}
\usepackage{comment}
\usepackage{minted}

% small et verbatim: to be coherent with \file
\newenvironment{smallverbatim}%
{\endgraf\small\verbatim}{\endverbatim}

\newenvironment{qv}
{\quote\Verbatim}
{\endquote\endVerbatim}

\usepackage{lastpage}

\newcommand{\terme}[1]{\texttt{#1}}
\newcommand{\termSet}[1]{\{\texttt{#1}\}}
\newcommand{\guill}[1]{<<\,#1\,>>}
% Macro pour l'inclusion de figure  
\newcommand{\fig}[2]{%
  \begin{figure*}[hbt]
   \begin{center}
      \includegraphics[width=6cm]{#1}
  \caption{\label{#1}  #2}
  \end{center}
 \end{figure*}
}
\newcommand{\figTex}[2]{%
  \begin{figure*}[hbt]
 \centerline{\input{#1.pstex_t}}
  \caption{\label{#1}  #2}
 \end{figure*}
}

% Macro commentaires
\definecolor{turquoise}{rgb}{0.20 ,0.60, 0.60}
\definecolor{vert}{rgb}{0.40,  0.60 ,0.00}
\definecolor{violet}{rgb}{0.80 , 0.00 , 0.80}
\definecolor{bleu}{rgb}{0.20,  0.20,  0.80}
\definecolor{rose}{rgb}{1.00,  0.20,  0.40}

\lstset{basicstyle=\footnotesize,showstringspaces=false,language=XML,breaklines=false,columns=fullflexible,frame=shadowbox, captionpos=b}

\newenvironment{solution}[0]{
~\\ \color{red} \textbf{Solution : }\color{black}}{%
}
\excludecomment{solution}


\newcommand{\cle}[1]{\textbf{#1}}
\def\fk#1{\textit{#1}}
\def\tble#1{\textit{#1}}

\parindent=0pt           
                             
\usepackage{url}

\newtheorem{Exercice}{Definition}


\begin{document}

\title{Sujet d'examen UE NFE204 - Session 1}
\author{Année universitaire 2022-2023}
\date{Examen première session: 31 janvier 2023}

\maketitle

\begin{center}
{\bf \Large Consignes}\\
\begin{tabular}{| p{15cm} |}
\hline
Dur\'ee de l'examen : 3h00; Documents autoris\'es\,: 10 pages recto-verso de notes manuscrites\,;  Calculatrice autoris\'ee.
Les exercices sont indépendants les uns des autres. Merci de soigner votre écriture.\\
Ce sujet contient \pageref{LastPage} pages.\\

\textit{Important}\,: nous n'évaluons pas les réponses par la syntaxe. Ne vous
inquiétez pas d'utiliser un point-virgule à la place d'un point d'exclamation ou
autres détails mineurs
\\
\hline
\end{tabular}
\end{center}


\section*{Première partie : documents structurés et MapReduce (7 pts)}

Nous gérons un site de commerce électronique qui recense des restaurants
et permet à des internautes de les noter. Pour l'instant nous utilisons
un système relationnel, dont voici le schéma (très simplifié).

\begin{itemize}
  \item Restaurant   (\textbf{id}, nom, catégorie)
  \item Internaute (\textbf{id}, nom)
  \item Notation (\textbf{\textit{idRestau}}, \textbf{\textit{idInternaute}}, note, commentaire)
\end{itemize}

\vspace*{0.2cm}

La catégorie du restaurant indique son niveau de prix, entre * et ****. 
Voici également un tout petit échantillon de la base.

\vspace*{0.2cm}
\begin{small}
\begin{minipage}[b]{4cm}
\begin{tabular}{l|l|l}
\cle{id} & nom & catégorie\\
\hline
20 & Bobo & ** \\
21 & Le régal & ****\\
\hline
\end{tabular}

\centerline{La table \tble{Restaurant}}
\end{minipage} \hspace*{0.5cm} \
\begin{minipage}[b]{2cm}
\begin{tabular}{l|l}
\cle{id} &  nom  \\
\hline
100 & Marie \\
101 & Adrien \\
\hline
\end{tabular}
\centerline{La table \tble{Internaute}}
\end{minipage} \hspace*{1cm} \
\begin{minipage}[b]{3cm}
\begin{tabular}{l|c|l|l}
\cle{\fk{idRestau}} & \cle{\fk{idInternaute}}  & note & commentaire \\
\hline
20  & 100  & 1 & Nul\\
20 & 101  & 5 & Parfait\\
21 & 101  & 3 & Moyen\\
\hline
\end{tabular}

\centerline{La table \tble{Notation}}
\end{minipage}
\end{small}

\vspace*{0.2cm}

On veut implanter  un algorithme de filtrage collaboratif. Un des calculs consiste notamment
à déterminer la note moyenne donnée par chaque internaute. 
Quelqu'un décide alors de remplacer la base relationnelle par une base NoSQL, et SQL par MapReduce.
Voyons s'il a raison.

\begin{enumerate}
  \item Donnez une représentation de l'échantillon ci-dessus sous la forme d'une collection
      de documents JSON, centrée sur les entités ``Restaurant''. 
      Nous l'appellerons représentation A.
  \item Proposez une autre représentation, centrée sur les entités ``Internaute''.  Nous l'appellerons représentation B.
             \item Donnez la requête SQL qui calcule la note moyenne donnée par chaque internaute.
           \item Expliquez le calcul MapReduce donnant cette moyenne, pour la représentation A.
           \item Même chose pour la représentation B.
           \item La phase de Reduce peut-elle être évitée pour l'un des calculs MapReduce\,?
                Si oui expliquez pourquoi, et l'impact sur la performance du calcul. 
\end{enumerate}

\textbf{Important}\,: pour spécifier les calculs MapReduce, donner les fonctions de Map et de Reduce, en javascript,
en pseudo-code, au pire en langage naturel en étant le plus précis possible.


\section*{Deuxième partie : recherche d'information (8 pts)}

On ajoute dans les  documents des commentaires sur les restaurants, désignés
par $c_1, ..., c_5$. On veut maintenant
créer un index plein texte pour rechercher des commentaires. On va travailler sur l'échantillon
suivant.

\vspace*{0.3cm}

\begin{tabular}{lp{15cm}}
     $c_1$& Plats médiocres, bonne carte des vins, accueil cordial.\\
     $c_2$& Le vin est moyen, les desserts sont mauvais, mais le décor est agréable.\\
     $c_3$& En général, le vin est trop cher au restaurant, mais ici la carte
         des vins est complète et de qualité.\\
     $c_4$& Passons sur le décor d'un goût douteux. Les plats et desserts sont délicieux,
          c'est tout ce qui compte après tout.\\
     $c_5$& Plat, dessert et un verre de vin pour un prix correct. Le plat du jour
             est toujours intéressant.\\ 
\end{tabular}

\bigskip

On va se limiter au vocabulaire (\texttt{vin, plat, décor, dessert}).
(NB: on suppose
     une phase préalable de normalisation qui élimine les pluriels, majuscules, etc.)

\begin{enumerate}
   \item Donnez une matrice d'incidence contenant les tf, 
      et un tableau donnant les idf (sans le $\log$), pour les  termes précédents. Placez
      les termes en ligne, les commentaires en colonne.

\begin{solution}
\begin{center}
\begin{tabular}{l|c|c|c|c|c}
            & \textbf{c1} & \textbf{c2} & \textbf{c3} &\textbf{c4}&\textbf{c5} \\ \hline  
   vin  5/4   & 1           & 1           & 2           & 0         & 1\\
   plat 5/3  & 1           & 0           & 0           & 1         & 2\\
   décor 5/2  & 0           & 1           & 0           & 1         & 0\\
   dessert 5/3  & 0           & 1           & 0           & 1         & 1\\\hline
   \end{tabular}
\end{center}
\end{solution}

  \item Donnez les normes des vecteurs représentant ces commentaires.
\begin{solution}
Les normes
\begin{itemize}
  \item $||c_1|| = \sqrt{1+ 1}$; $||c_2|| = \sqrt{3}$; $||c_3|| = \sqrt{4}$; 
    $||c_4|| = \sqrt{3}$, $||c_5|| = \sqrt{6}$  
\end{itemize}
\end{solution}
  
  \item Donner les résultats classés par similarité cosinus basée sur les tf 
  (on ignore l'idf) pour les requêtes suivantes. 
  
  \begin{itemize}
    \item plat;
    \item décor et dessert;
    \item vin et plat.
  \end{itemize}

Expliquez brièvement la raison  du classement pour le premier document.

\begin{solution}
Les cosinus (requête non normalisée).

\begin{itemize}
  \item plat\,: c1 : $\frac{1}{\sqrt{2}}$;  
       c4 : $\frac{1}{\sqrt{3}}$; 
       c5 : $\frac{2}{\sqrt{6}}$; 
  
  \item décor et dessert\,: 
  
    c2 : $\frac{1+1}{\sqrt{3}}$; 
    c4 :  $\frac{1+1}{\sqrt{3}}$  
    c5 : $\frac{1}{\sqrt{6}}$; 

  \item vin et plat
  
    c1 : $\frac{1+1}{\sqrt{2}}$; 
    c2 : $\frac{1}{\sqrt{3}}$;
    c3 :  $\frac{2}{\sqrt{3}}$  
    c4 :  $\frac{1}{\sqrt{3}}$  
    c5 : $\frac{1+2}{\sqrt{6}}$; 
    
 \end{itemize}

\end{solution}  

   \item Pour la requête "décor et dessert", et les documents c2 et c4,
   qu'est-ce qui change si on ne met pas de restriction 
   sur le vocabulaire (tous les mots sont indexés)?
   
\begin{solution}
Dans ce cas le plus petit document (c2) l'emporte car sa norme est
plus petite. Intuititvement: il parle plus \emph{spécifiquement} de
décor et de dessert que c4.
\end{solution}   
    
\item On voudrait pouvoir mesurer les co-occurrences de termes, autrement dit
    produire un indicateur représentant la probabilité que deux termes
    apparaissent ensembles dans un même commentaire. Proposez une formule estimant
    cette probablilité et s'appuyant
      sur les informations présentes dans la matrice d'incidence.

\end{enumerate}


\section*{Troisième partie : flux (2 pts)}

Nous recevons un flux de notes données à des restaurants nommé \texttt{restaus}. Chaque
élément contient le nom du restaurant, l'identifiant
de l'utilisateur, la note et le commentaire. Voici un échantillon.

\begin{smallverbatim}
1,Le bougnat,100,4,"Plats médiocres ..."
2,Délices de Saïgon,100,2,"Desserts nuls..."
3,Le régal,101,4,"venez pour le décor"
..
\end{smallverbatim}

On met en place une chaîne de traitement pour ce flux. Voici cette chaîne exprimée avec les opérateurs de Pig 

\begin{smallverbatim}
coll1 = filter restaus by contains(commentaire, "dessert")
coll2 = group coll1 by nomRestaurant;
coll3  = foreach coll2 generate group, AVG(coll1.note);
\end{smallverbatim} 

\begin{enumerate}
 \item Que produit cette chaîne\,?
 

\begin{solution}
La note moyenne, par restaurant, des commentaires relatifs aux desserts
\end{solution}
 
 \item Quel est le degré maximal de parallélisation\,?
\begin{solution}
La parallélisation est limitée par le nombre de restaurants.
\end{solution}

\end{enumerate}

\section*{Quatrième partie : questions de cours (3 pts)}

\begin{enumerate}
\item Je soumets une requête $t_1, t_2, \cdots, t_n$. 
    Quel est le poids de chaque terme dans le vecteur représentant cette requête? 
    La normalisation de ce vecteur est elle importante pour le classement (justifier)?
  
    \item Dans un traitement MapReduce, à quoi correspond l'opération dite de \textit{shuffle}?
    
    \item Dans un système distribué avec partitionnement, je
          décide de partitionner ma collection de restaurants
          sur la catégorie. Discuter des avantages et inconvénients de ce choix.   
\end{enumerate}

\end{document}
